\documentclass[11pt,a4paper]{article}

\usepackage[utf8]{inputenc}
\usepackage[T1]{fontenc}
\usepackage[french]{babel}
\usepackage[margin=2.5cm,headheight=14pt]{geometry}
\usepackage{lmodern}
\usepackage{tikz}
\usetikzlibrary{arrows.meta,positioning,fit,babel}
\usepackage{listings}
\usepackage{booktabs}
\usepackage{float}
\usepackage{fancyhdr}
\usepackage[hidelinks]{hyperref}

\addto\captionsfrench{\renewcommand{\tablename}{Tableau}}

% ---------- En-tête et pied de page ----------
\pagestyle{fancy}
\fancyhf{}
\fancyhead[L]{\small TP1 -- Virtualisation}
\fancyhead[R]{\small ERRAMI YOUSSEF}
\fancyfoot[C]{\thepage}

% ---------- Blocs de code ----------
\renewcommand{\lstlistingname}{Listing}
\lstset{
  basicstyle=\ttfamily\small,
  frame=single,
  framesep=6pt,
  captionpos=b,
  breaklines=true,
  showstringspaces=false,
  columns=fullflexible,
  keepspaces=true
}

% ---------- Styles des figures ----------
\tikzset{
  boite/.style={draw, minimum width=2.5cm, minimum height=1.1cm,
    align=center, font=\small},
  fleche/.style={-{Stealth[length=2mm]}},
  retour/.style={fleche, dashed},
  etiq/.style={font=\scriptsize, align=center}
}

\title{Déploiement d'une application multi-conteneurs\\avec Docker Compose}
\author{ERRAMI YOUSSEF\\[2pt] \small TP1 -- Virtualisation}
\date{\today}

\begin{document}
\maketitle
\thispagestyle{fancy}

\section{Introduction}
L'objectif de ce travail pratique est de réaliser une application qui affiche les
dix premières chansons du classement iTunes. L'application est répartie en deux
conteneurs ayant chacun un rôle précis :
\begin{itemize}
  \item \texttt{fetcher} : récupère le flux XML d'Apple, en extrait les informations
        utiles et les convertit au format JSON ;
  \item \texttt{app} : interroge le service \texttt{fetcher} et affiche les chansons
        dans une page web.
\end{itemize}

La source des données est l'adresse suivante :\\
{\small\url{http://ax.itunes.apple.com/WebObjects/MZStoreServices.woa/ws/RSS/topsongs/limit=10/xml}}

\section{Architecture}
La figure~\ref{fig:architecture} présente le cheminement d'une requête depuis le
navigateur jusqu'au serveur d'Apple, ainsi que le retour de la réponse.

\begin{figure}[H]
\centering
\begin{tikzpicture}
  \node[boite] (nav) at (0,0)    {Navigateur};
  \node[boite] (app) at (4.3,0)  {\texttt{app}\\port 3000};
  \node[boite] (fet) at (8.7,0)  {\texttt{fetcher}\\port 4000};
  \node[boite] (apl) at (13,0)   {Apple\\flux XML};

  \draw[fleche] ([yshift=2.5mm]nav.east) -- node[etiq, above]{GET /} ([yshift=2.5mm]app.west);
  \draw[retour] ([yshift=-2.5mm]app.west) -- node[etiq, below]{HTML} ([yshift=-2.5mm]nav.east);
  \draw[fleche] ([yshift=2.5mm]app.east) -- node[etiq, above]{/songs} ([yshift=2.5mm]fet.west);
  \draw[retour] ([yshift=-2.5mm]fet.west) -- node[etiq, below]{JSON} ([yshift=-2.5mm]app.east);
  \draw[fleche] ([yshift=2.5mm]fet.east) -- node[etiq, above]{requête} ([yshift=2.5mm]apl.west);
  \draw[retour] ([yshift=-2.5mm]apl.west) -- node[etiq, below]{XML} ([yshift=-2.5mm]fet.east);

  \node[draw, dotted, thick, fit=(app)(fet), inner sep=4mm,
    label={[font=\small]above:Réseau Docker Compose}] {};
\end{tikzpicture}
\caption{Architecture de l'application}
\label{fig:architecture}
\end{figure}

\section{Dockerfile}
Un Dockerfile contient les instructions de construction d'une image. Docker les
exécute de haut en bas. Le listing~\ref{lst:dockerfile} correspond au service
\texttt{app}.

\begin{lstlisting}[caption={Fichier app/Dockerfile}, label={lst:dockerfile}]
FROM node:22-alpine
WORKDIR /app

# Dependances d'abord pour profiter du cache Docker
COPY package.json ./
RUN npm install --omit=dev

COPY server.js ./
USER node
EXPOSE 3000
CMD ["node", "server.js"]
\end{lstlisting}

\begin{table}[H]
\centering
\begin{tabular}{lp{11.5cm}}
\toprule
\textbf{Instruction} & \textbf{Rôle} \\
\midrule
\texttt{FROM}    & Image de base : Node.js 22 sur la distribution légère Alpine Linux. \\
\texttt{WORKDIR} & Définit le répertoire de travail dans l'image. \\
\texttt{COPY}    & Copie des fichiers de la machine hôte vers l'image. \\
\texttt{RUN}     & Installe les dépendances pendant la construction. \\
\texttt{USER}    & Exécute l'application sans les droits administrateur. \\
\texttt{EXPOSE}  & Documente le port utilisé, sans le publier sur la machine hôte. \\
\texttt{CMD}     & Définit la commande exécutée au démarrage du conteneur. \\
\bottomrule
\end{tabular}
\caption{Instructions du Dockerfile}
\end{table}

\subsection{Les couches de l'image}
La figure~\ref{fig:couches} montre l'empilement des principales étapes de
construction. Le fichier \texttt{package.json} est copié avant le code source :
lorsque seul \texttt{server.js} est modifié, Docker réutilise le cache et ne
réinstalle pas les dépendances.

\begin{figure}[H]
\centering
\begin{tikzpicture}[couche/.style={draw, minimum width=7cm, minimum height=6.5mm,
    font=\small\ttfamily}]
  \node[couche] (a) at (0,0)    {FROM node:22-alpine};
  \node[couche] (b) at (0,0.65) {WORKDIR /app};
  \node[couche] (c) at (0,1.3)  {COPY package.json};
  \node[couche] (d) at (0,1.95) {RUN npm install};
  \node[couche] (e) at (0,2.6)  {COPY server.js};
  \draw[fleche] (-4,-0.3) -- node[etiq, left]{ordre de\\construction} (-4,2.9);
  \node[etiq, right=3mm of e] {modifiée souvent};
  \node[etiq, right=3mm of c] {modifiée rarement};
\end{tikzpicture}
\caption{Couches de l'image du service \texttt{app}}
\label{fig:couches}
\end{figure}

Le Dockerfile du service \texttt{fetcher} suit le même modèle, avec l'instruction
\texttt{EXPOSE 4000}. Dans les deux dossiers, un fichier \texttt{.dockerignore}
exclut \texttt{node\_modules} du contexte de construction.

\section{Docker Compose}
Docker Compose permet de décrire et de lancer plusieurs conteneurs à partir d'un
seul fichier (listing~\ref{lst:compose}).

\begin{lstlisting}[caption={Fichier docker-compose.yml}, label={lst:compose}]
services:
  # Va chercher les chansons chez Apple (aucun port publie)
  fetcher:
    build: ./fetcher

  # Affiche les chansons
  app:
    build: ./app
    ports:
      - "3000:3000"
    environment:
      FETCHER_URL: http://fetcher:4000
    depends_on:
      - fetcher
\end{lstlisting}

\begin{table}[H]
\centering
\begin{tabular}{lp{11.5cm}}
\toprule
\textbf{Clé} & \textbf{Rôle} \\
\midrule
\texttt{services}    & Liste des conteneurs de l'application. \\
\texttt{build}       & Dossier contenant le Dockerfile à construire. \\
\texttt{ports}       & Relie le port 3000 de l'hôte à celui du conteneur. \\
\texttt{environment} & Transmet à \texttt{app} l'adresse du service \texttt{fetcher}. \\
\texttt{depends\_on} & Démarre \texttt{fetcher} avant \texttt{app}, sans attendre qu'il soit prêt. \\
\bottomrule
\end{tabular}
\caption{Éléments du fichier Docker Compose}
\end{table}

\subsection{Communication entre les conteneurs}
Compose crée un réseau commun dans lequel chaque service est joignable par son nom.
L'application appelle donc \texttt{http://fetcher:4000/songs} sans connaître
l'adresse IP du conteneur. Le port 4000 n'étant pas publié, seule l'application est
accessible depuis la machine hôte.

\section{Exécution}
Les commandes, définies dans le fichier \texttt{package.json} à la racine du projet,
sont résumées dans le tableau~\ref{tab:commandes}.

\begin{table}[H]
\centering
\begin{tabular}{lll}
\toprule
\textbf{Commande} & \textbf{Commande Docker} & \textbf{Rôle} \\
\midrule
\texttt{npm run dev}  & \texttt{docker compose up --build} & Construire et démarrer \\
\texttt{npm run logs} & \texttt{docker compose logs -f}    & Suivre les journaux \\
\texttt{npm run stop} & \texttt{docker compose down}       & Arrêter et supprimer \\
\bottomrule
\end{tabular}
\caption{Commandes d'exécution}
\label{tab:commandes}
\end{table}

Après le démarrage, la page web est disponible à l'adresse
\url{http://localhost:3000} et les données JSON à l'adresse
\url{http://localhost:3000/api/songs}.

\section{Conclusion}
Les deux conteneurs ont été construits et lancés avec Docker Compose, et
l'application a renvoyé les dix chansons attendues. Ce travail pratique a permis
d'appliquer l'écriture d'un Dockerfile, l'organisation de plusieurs services avec
Compose et la communication entre conteneurs. La séparation des rôles permet de
changer la source des données sans modifier l'affichage.

\end{document}
